\subsection{Early Exit Networks} % 改叫early exit networks?.
The early exit network is designed to reduce the computational resource during the model inference time while enhancing the model accuracy. With the early exit scheme, early exit branches will be attached to DNN. When the inference phase, the instances will stop inferencing and leave the model from different branches following the pre-defined early exit rules or criteria. BranchyNet \cite{Teerapittayanon2016BranchyNetFI}  calculates entropy at each branch after obtaining the result and exits if the entropy of the predicted result is less than the threshold value. Shallow-Deep Networks \cite{Kaya2019ShallowDeepNU} calculates the softmax score of each internal classifier's prediction and takes the maximum probability value as the confidence score. Once the score exceeds the threshold during the forward passing, the instance will exit from the branch prematurely, and the model further mitigates the "overthinking" problem of DNNs. \cite{schwartz-etal-2020-right} leveraged the early exit inference scheme of \cite{Kaya2019ShallowDeepNU} and applied a confidence-based strategy to the NLP task. \cite{Zhou2020BERTLP} makes predictions using adjacent layers and stops inference when the predicted value of the internal classifier remains constant in a given inference unit times. We borrow the architecture of \cite{Kaya2019ShallowDeepNU} and the confidence-based early exit strategy to achieve fairness classification. To our best knowledge, we are the first work that leverages the early exit technique to improve fairness.


\subsection{Bias Mitigation Methods}
Bias mitigation methods are designed to reduce the native bias in the dataset to reduce the chance of unfair prediction. There are largely three avenues for current debiased strategy, including pre-processing, in-processing, and post-processing. 
(1) Pre-processing strategies usually remove the information which may cause "discrimination" from training data before training. \cite{Kamiran2011DataPT} use different weight to neutralize the effect of the sensitive information in training phase and present the experiments result on real-life data. \cite{Ngxande2020BiasRI} and \cite{Lu2020GenderBI} achieve fairness via data pre-process methods including data generation and data augmentation.

In-processing method usually modify on-the-shelf model architecture, loss function, and model regularization to achieve fairness goal. \cite{Wang2020TowardsFI} used domain-independent training technique to improve fairness. Adversarial training mitigates the bias through adversarially trains an encoder and classifier to learn a fair representation. \cite{Zhang2018MitigatingUB} adversarially cooperate a predictor and a adversary to make accurate prediction while removing the sensitive attributes from the representation. \cite{kim2019learning} unlearned the bias of the dataset with the proposed regularization and trained the model in a adversarial fashion.
Disentangled representation forces the encoder abandon the information of sensitive attributes and thus achieve fairness.  

Post-processing method aims to calibrate the model's output to enhance fairness. They need to use sensitive attribute and prediction distribution to modify the previous distribution result. The \cite{Hardt2016EqualityOO} reveals the limit of demographic parity and give a new metric to fairness, equalized odds, and show how to adjust the learned prediction. \cite{Zhao2017MenAL} used Lagrangian relaxation to designed an inference algorithm which decrease bias but maintain performance in the meantime.


%% ?????
% However, those pre-processing and post-processing methods are not suitable for medical image diagnosis task because they need to remove or access sensitive attribute during experiment, it may influence the diagnosis, so our method focus on improving fairness in training phase. Moreover, out method can achieve fairness without knowing sensitive attribute in the training time, which is a "fairness through unawareness" technique.

However, most pre-processing and post-processing methods are unsuitable for medical image diagnosis tasks because they are required to remove or access sensitive attributes during the experiment, which are not practical. As a result, our method focus on improving fairness in the training phase. Moreover, our method can achieve fairness without knowing sensitive attributes during training, which is a "fairness through unawareness" technique.